\chapter{Sintonización PID mediante Eagle Strategy}
\label{ch:eagle-strategy}

\section{Introducción del capítulo}
Eagle Strategy no suele presentarse como un algoritmo único y cerrado, sino como un marco de
búsqueda que alterna exploración global e intensificación local. Su inclusión en esta obra
resulta útil porque permite discutir una idea muy relevante en optimización aplicada: no todos
los problemas se benefician de una única dinámica de movimiento; a veces conviene recorrer
rápidamente regiones amplias y luego afinar solo donde el paisaje de soluciones lo justifica.

\section{Objetivo del capítulo}
El objetivo de este capítulo es exponer la lógica general de Eagle Strategy y adaptarla a la
sintonización de controladores PID. Más que profundizar en variantes matemáticas complejas,
se busca mostrar cómo una estrategia híbrida de dos fases puede organizar la búsqueda de
ganancias cuando la función objetivo presenta múltiples regiones de interés.

\section{Fundamento del método}
La idea básica del método consiste en alternar dos momentos. En la primera fase, el algoritmo
realiza desplazamientos amplios, frecuentemente asociados a saltos aleatorios de tipo Levy, con
el fin de descubrir nuevas regiones prometedoras del espacio de búsqueda. En la segunda
fase, una técnica local o semilocal refina la solución encontrada en ese vecindario. El nombre
del método alude precisamente a esa conducta: detectar desde lejos y atacar con precisión.

\section{Estructura algorítmica}
En una aplicación al ajuste PID, la fase global puede entenderse como una exploración sobre
rangos amplios de $K_p$, $K_i$ y $K_d$, mientras que la fase local se concentra en pequeñas
perturbaciones alrededor de las mejores configuraciones observadas. Esta estructura reduce el
riesgo de estancamiento temprano y permite aprovechar información local sin renunciar a la
capacidad de descubrir soluciones alejadas de la inicialización.

\section{Pseudocódigo}
\begin{algorithm}[H]
\caption{Esquema general de Eagle Strategy para sintonización PID}
\begin{algorithmic}[1]
\State Inicializar una solución o conjunto reducido de soluciones
\While{no se cumpla el criterio de parada}
    \State Ejecutar una fase de exploración global sobre el espacio $\left[K_p,K_i,K_d\right]$
    \State Seleccionar la región o solución más prometedora
    \State Aplicar una búsqueda local para refinar la sintonización
    \State Actualizar la mejor solución conocida
\EndWhile
\State Entregar la tripleta de ganancias con menor costo
\end{algorithmic}
\end{algorithm}

\section{Parámetros relevantes}
Los parámetros relevantes de esta estrategia dependen de la implementación elegida, pero
normalmente incluyen amplitud de los saltos globales, tamaño del vecindario local, número de
refinamientos por ciclo y criterio de cambio entre fases. Esta flexibilidad puede ser una
fortaleza porque permite adaptar el método a distintos problemas; sin embargo, también exige
una calibración más cuidadosa que la requerida por algoritmos con estructura más estandarizada.

\section{Aplicación a la sintonización PID}
Aplicada al problema de sintonización, Eagle Strategy puede utilizarse para detectar primero
zonas del espacio de ganancias con buen compromiso entre rapidez y amortiguamiento, y luego
afinar la respuesta en torno a esas zonas usando perturbaciones controladas. Esto la vuelve
interesante para problemas donde la evaluación del ITAE muestra varios valles locales o donde
una exploración exclusivamente intensiva podría perder soluciones útiles. En la práctica, su
desempeño debe juzgarse con el mismo marco comparativo usado para PSO y GWO.

\section{Cierre del capítulo}
Eagle Strategy aporta una lectura híbrida del problema y refuerza la idea de que una buena
sintonización depende tanto de la función objetivo como de la manera en que se recorre el
espacio de búsqueda. El capítulo siguiente cierra el bloque de métodos con Artificial Bee
Colony, otra propuesta colectiva inspirada en mecanismos cooperativos de búsqueda.
